\documentclass[11pt]{article}
\usepackage{setspace}
\usepackage{geometry}
\geometry{margin=1in}
\setstretch{1.15}

\begin{document}

\noindent
Zhongli Wang \\
Bielefeld, Germany \\
Phone: +49 160 9541 6015 \\
Email: nkwangzhongli@gmail.com \\
Website: https://sites.google.com/view/zhongli-wang \\

\vspace{0.35cm}

\noindent
Hiring Committee \\
Fideres \\
Competition Litigation Consulting \\

\vspace{0.35cm}

\noindent
\textbf{Application for Competition Economist} \\

\vspace{0.25cm}

\noindent Dear Hiring Committee, \\[0.15cm]

\noindent
I am writing to apply for the Competition Economist position at Fideres (I'm open to Madrid or London). I am completing my Ph.D.\ in Economics at Bielefeld University and Universit\`a Cattolica del Sacro Cuore (Marie Sk{\l}odowska--Curie EPOC programme). My research sits at the intersection of industrial organization, applied econometrics, and computational methods, with a focus on how strategic incentives and information environments shape market outcomes. Fideres' litigation-only model and emphasis on rigorous, conflict-free analysis strongly appeal to me: I am motivated by work where careful empirical reasoning, transparent documentation, and clear communication can directly influence high-stakes decisions.

\vspace{0.2cm}

\noindent
My academic work has trained me to do exactly what competition litigation requires: (i) diagnose the core economic mechanism in a complex fact pattern, (ii) translate it into a testable framework, and (iii) produce defensible, well-documented results under tight deadlines. In my dissertation, \textit{Strategic Behavior and Economic Dynamics in ICT Markets: An Agent-Based Framework}, I study competition in technology-intensive markets characterized by network effects, congestion, and switching frictions. I build and estimate dynamic models in which firms adjust pricing and investment in response to demand feedback and rivals' strategies, generating disciplined counterfactuals about market power, consumer harm, and the effects of alternative market rules. This work has sharpened my ability to reason about platform competition, barriers to entry, and the link between market design and welfare---topics that frequently arise in modern competition disputes.

\vspace{0.2cm}

\noindent
In parallel, I have pursued research on financial market behavior and transparency that complements litigation settings where incentives and information manipulation are central. In \textit{Polytope Fraud Theory}, I use supervised machine learning to detect structured irregularities in firm financial statements, focusing on firms targeted by activist short sellers in U.S.\ markets. The project combines careful data construction, feature engineering, and robustness checks with an economic interpretation of incentives, monitoring, and strategic misreporting. The result is a practical framework that connects observable anomalies to institutional mechanisms---a workflow that is directly relevant to securities and market-manipulation contexts, and to competition matters where the credibility of disclosures, internal documents, or claimed efficiencies must be assessed empirically.
In addition, my coauthored paper paper \textit{Agentic Workflows for Economic Research} introduces a general framework for constructing LLM-based agentic workflows, including retrieval-augmented generation, prompt engineering, evaluation metrics, and pipeline governance. 
This project demonstrates my capacity to design, implement, and audit advanced AI systems—skills that directly complement emerging work in NLP-driven analysis, causal inference using unstructured data, and the development of AI tools for field experiments.

\vspace{0.2cm}

\noindent
Beyond topic fit, I would bring to Fideres a disciplined and reproducible analytical style. I work primarily in Python (and can work in R when required), and I routinely build end-to-end pipelines: data ingestion and cleaning, exploratory analysis, econometric estimation, sensitivity/robustness checks, and production of client-ready tables and figures. I am comfortable working with large datasets and integrating multiple sources, and I have experience with high-performance computing for simulation and computationally intensive estimation. Importantly, I have learned to communicate technical results in plain language---including assumptions, limitations, and the precise interpretation of estimates---so that non-specialists can understand what the analysis does and does not establish. I am also fluent with Excel, Word, and PowerPoint for producing concise, well-structured materials tailored to legal and business audiences.

\vspace{0.2cm}

\noindent
The aspects of the role that excite me most are the ones where Fideres is distinctive: developing case theories into concrete empirical strategies, iterating quickly as new facts emerge, and drafting clear reports that withstand scrutiny. I enjoy working in teams, taking ownership of deliverables, and proactively improving processes (e.g., reusable code modules, standardized QC checks, and templates for transparent reporting). I would welcome the opportunity to contribute to projects involving Big Tech and other complex markets, and to support senior colleagues in preparing analyses that are both technically sound and persuasive in their presentation.

\vspace{0.2cm}

\noindent
I would be delighted to discuss how my background in industrial organization, empirical analysis, and computational modelling can support Fideres' competition litigation practice. Thank you so much for your consideration!

\vspace{0.45cm}

\noindent
Sincerely,\\[1cm]
Zhongli Wang

\end{document}
